% ===== §10 The Generalized Field of Travel =====
\section{The Generalized Field of Travel}
\label{p21:sec:generalized}

\noindent
The paper has spoken throughout of travel in its literal sense, the journey to an unfamiliar place. It is time to generalise. ``Travel,'' it will be argued, names not essentially a movement through geographical space but a \emph{function}: the construction of the special field, the opening of the spacious, contingent, suspended condition within which a relation comes to presence. The literal journey is the most efficient and powerful instance of this function, but it is an instance, not the essence; and once the function is distinguished from its most familiar instance, it becomes possible to perform it without the journey, to construct the field within the everyday, with neither distance nor expense. This generalisation is the practical culmination of the paper, and it is offered not as a conclusion that closes the argument but as an opening that hands it on. It is also the resolution of a worry that has shadowed the paper from the start: that its goods are the privileges of those who can afford to travel.

% ============================================================
\subsection{Travel as a functional concept}
\label{p21:sub:gen-functional}

What, in the end, did the literal journey do? It opened the special field: it exchanged the fixed background, suspended the roles, restored contingency, framed a bounded continuous time, threw two people into symmetric unfamiliarity and forced co-presence, and beneath all these, opened the spaciousness in which a relation could come to presence. These are the functions; the journey is their vehicle. And once the functions are named in their own right, it becomes clear that they are not necessarily tied to geographical movement. The exchange of background, the suspension of roles, the restoration of contingency, the opening of spaciousness, each of these can, in principle, be effected without going anywhere, by other means than the journey. The journey effects them all at once, powerfully and reliably, which is its great advantage; but it holds no monopoly on them. To generalise travel is to recognise that ``travel,'' in the sense that has mattered to this paper, \emph{is} the performance of these functions, the construction of the special field, and that the literal journey is one way, the most efficient way, but not the only way, of performing it.

This recognition reorganises the whole paper retrospectively. Every feature of the field analysed in \xref{p21:sec:field} is, seen now, a function that can be sought in the everyday: one can exchange the background by altering the familiar setting; one can suspend roles by deliberately and jointly bracketing the everyday's demands for a protected span; one can restore contingency by introducing the undetermined into the over-determined day; one can open spaciousness by refusing, for a while, the everyday's total assignment of function. The generalised field of travel is the everyday deliberately opened into a field, the construction, without the journey, of the spacious contingent condition the journey opens with such ease. And the deepest of the functions, the one that underlies all the others and is the most portable of all, is the one the aesthetic chapters identified: the perception that catches the contingent and the spaciousness that lets it be caught. For if the contingent is always already leaking through the seams of the everyday (\xref{p21:sub:ap-everywhere}), then the field is, in a sense, always already half-open, awaiting only the perception that would catch what it offers. To generalise travel is, at its root, to carry the perception home.

% ============================================================
\subsection{From scaffolding to skill}
\label{p21:sub:gen-scaffolding}

The relation between the literal journey and the generalised field is the relation, introduced earlier, between scaffolding and skill (\xref{p21:sub:ap-skill}). The literal journey is a scaffolding: a powerful external structure, erected at some cost, that opens the field reliably and completely, holding it open by main force of distance, expense, and the wholesale suspension a journey effects. Its value is double. It opens the field, here and now, for the two who travel; and, the deeper value, it \emph{teaches the field}, lets two people experience the special field vividly enough to learn what it is, to recognise its character, to know from the inside what it is to be in it and to catch and share within it. The journey is the demonstration class in which one learns, by undergoing it, what the field is and what one does within it.

But the mature form of the practice is not the perpetual erection of scaffoldings. To require a literal journey every time one wished to open the field would be to remain forever dependent on an expensive external structure, and it is exactly this dependence that would make the paper's goods the privilege of the moneyed and the free. The mature form is the \emph{skill} that the scaffolding taught: the internalised capacity to open the field, or to catch what the half-open everyday already offers, without the external structure. Once two people have learned the field, learned it, perhaps, through literal journeys that showed them what it is, they can begin to open it at home, in the ordinary day, by the cultivated capacities the journeys taught them: the perception that catches the contingent, the taste that constructs the frame, the restraint that keeps the vessel empty, the mutual attention that turns the caught beauty toward the other. The scaffolding comes down; the skill remains; and the skill opens, in the midst of the everyday, the field the scaffolding was erected to open. This is the movement the whole paper has been tending toward: from the journey that opens the field by force, to the skill that opens it anywhere; from travel as a thing one does occasionally and expensively, to travel as a capacity one carries always and exercises freely. The literal journey is the demonstration; the generalised field is the daily practice; and the bridge between them is the cultivation, through the one, of the skill that performs the other.

% ============================================================
\subsection{The de-alienation of the everyday, and the answer to elitism}
\label{p21:sub:gen-dealienation}

This resolves the worry of elitism, and the resolution is also the paper's deepest political claim. The worry is obvious and serious: if the good this paper describes required literal travel, it would be the privilege of those with the money and the leisure to travel, and a philosophy of intimacy resting on it would be a philosophy for the comfortable, complicit in the very inequality the series has elsewhere opposed \parencite{paperxv,paperxvii}. The generalisation dissolves the worry. The generalised field does not require distance or expense; it requires the cultivated capacities, perception, taste, restraint, attention, which depend on no wealth and are, as the previous section argued, exactly the capacities that capital can use but cannot confiscate and that the poorest perceiver may possess more fully than the richest tourist (\xref{p21:sub:pe-spectrum}). The good of the field is therefore, in its generalised form, \emph{universally available}: open to any two people willing to cultivate the perception and do the relational labour, regardless of their means. The literal journey is a privilege; the generalised field is not, and the movement from the one to the other is the movement from a good available to the few to a good available to all.

This is, in the terms the previous section established, the \emph{de-alienation of the everyday}. The alienated everyday is the everyday wholly assigned to function, every hour given to labour or its reproduction, every space to use, the whole of life instrumentalised under the administered order, with no opening for the non-instrumental, the gratuitous, the spacious \parencite{marx1992,debord1994}. To open a field within this everyday, to clear a protected span in which roles are suspended and the purposive lifted, to catch the contingent beauty leaking through the seams, to undergo a moment of genuine non-instrumental experience with the beloved, is to reclaim a piece of the everyday from its alienation, to make, within the administered order, a small space that is genuinely one's own. And because this reclamation requires no wealth, it is available where wealth is not: the de-alienation of the everyday is a freedom the poor can practice as fully as the rich, perhaps more fully, since the cultivated perception that is its only requirement is no respecter of means. Here the paper's politics reaches its conclusion. The cultivation of the field's skill is the building, by two people, of a capacity to de-alienate their everyday, to open, within the administered world and without its permission, recurring small fields of genuine, spacious, non-instrumental, shared experience. This is not a retreat from the political into the domestic; it is a politics conducted at the scale of the everyday, the construction of a commons of de-alienated experience that the administered order cannot enclose because it requires nothing the order controls. The generalised field of travel is the everyday reclaimed, and the reclaiming is available to all.

% ============================================================
% ============================================================
\subsection{Generative openness: a principle, and a limit}
\label{p21:sub:gen-openness}

The generalisation of travel permits one further abstraction, which the paper states as a principle and then, characteristically, marks as opening beyond itself. Throughout, the condition underlying the field has been spaciousness (\xref{p21:sub:field-spaciousness}), the emptiness in the sense of the \emph{Daodejing}'s eleventh chapter, the room left unfilled in which the contingent can arrive and be caught. The generalisation now lets this be seen not as a feature peculiar to travel but as a general \emph{resource}, the openness that any field requires if relational generativity is to occur within it, and the resource can be named: \emph{generative openness}, the condition of not being so wholly filled by function that nothing unplanned can enter. A field, of whatever kind, is generative to the degree that it preserves this openness; and the principle may be stated generally: \emph{a field is by preserving the openness in which unplanned, meaningful events may emerge} rather than made generative by adding activity to it,. The political economy of beauty already confirmed this from the side of value (\xref{p21:sub:pe-emptiness}): the costly emptiness that capital pays dearly to preserve is generative openness recognised, even by capital, as a resource rather than a waste.

Two qualifications keep the principle from being misunderstood, and the second of them marks the paper's limit. The first: generative openness is not emptiness as a style. It must not be confused with minimalism, or with an aesthetic of bareness, for the point is that room must be left for the unplanned rather than that less is better, and that room may be held open as well by a thing that invites lingering (a window, a bench, a painting, a piano, a table wide enough to gather around) as by literal vacancy. Generative openness is the condition of not being wholly closed by function, whatever holds that condition open; it is a structural property of a field, not a decorative quantity of empty space. The second qualification is the limit. The principle of generative openness plainly reaches beyond the intimate relation that is this paper's subject, it bears on the design of rooms and houses, of schools and libraries and cities, of institutions and their allocation of time and budget and attention, wherever the question arises whether a space or a structure preserves the openness in which generative events can occur or forecloses it by total functionalisation. But the full development of generative openness as a principle of design, spatial, temporal, institutional, is not the work of this paper, and would, if pursued here, carry the argument away from the relational field it is meant to illuminate into a general theory of generative design the paper cannot contain. The paper therefore states the principle, notes that the field of travel and the costly emptiness of the gallery and the unfilled room are instances of it, and marks its general development as a directed opening (\xref{p21:sub:clo-fissures}), a principle this paper discovers at the limit of its own subject and hands on to a theory of generative design it borders but does not enter. That generative openness is a real and general resource, the paper affirms; what a full theory of its cultivation across spaces, times, and institutions would be, it leaves, deliberately, to other work.

% ============================================================
\subsection{An opening, not a conclusion}
\label{p21:sub:gen-opening}

It would be false to the matter to present the generalised field as a solution achieved, a technique mastered, a conclusion reached. The generalisation opens a practice; it does not complete one. To open a field within the everyday is difficult, and to do it repeatedly, sustainably, against the everyday's constant pressure to close back down, is more difficult still, a labour without end, a practice that must be perpetually renewed because the everyday perpetually re-alienates, the opened field perpetually re-closing into routine. The paper does not pretend to have mastered this practice or to be able to hand the reader a method for it; \el{χαλεπὰ τὰ καλά}, and the generalised field, being an aesthetic accomplishment, cannot be reduced to a method any more than the good or the beautiful can. What the paper offers is the recognition that the practice exists, that its direction is the movement from scaffolding to skill, and that its goal is the daily, mutual, un-programmable cultivation of the capacity to open fields of de-alienated shared experience in the midst of an everyday that is forever closing them. How to sustain this practice, how to keep the field open against the everyday's pressure, how not to let the skill decay back into the routine it was meant to interrupt, these are real and open questions, and the paper hands them on rather than answering them. The generalised field is not the paper's conclusion but its opening: the point at which the understanding it has built passes into a practice that exceeds it, and the reader is left not with a doctrine but with a direction and a difficulty. To which the only fitting response is the one the whole paper has been arguing for: not a method, but the cultivation, slow and mutual and lifelong, of the perception and the taste and the love that open the field, and the difficult, unending, beautiful labour of keeping it open, together.
